% ============================================================================
% EXAMPLE: Single-paper talk (auto-generated)
%
% Produced by running:
%   python scripts/plan_deck.py /tmp/bundle_test \
%     --output-dir /tmp/final_deck \
%     --title "How Memory Works: A Retrieved-Context Account" \
%     --subtitle "Talk based on our 2024 manuscript" \
%     --target-slides 15
%
% on a fake paper bundle. Preserved here as a reference for what the
% skill emits before deep-reading fills in the [TODO ...] placeholders.
% Compiles with pdflatex + bibtex to an 11-slide PDF.
% ============================================================================
% ============================================================================
% CML Beamer Template
% Matches the visual style used in Computational Memory Lab lectures
% (Madrid theme + navy blue + lab-branded footer).
%
% Variables to fill in (look for <<...>> markers):
%   How Memory Works: A Retrieved-Context Account       main talk title
%   Talk based on our 2024 manuscript    optional subtitle (delete the \subtitle line if unused)
%   Alice Smith, Bob Jones      presenter name(s)
%   \today        defaults to today
%   How Memory Works: A Retrieved… short version shown in footer (max ~5 words)
%
% Use \section{...} for each major section. Each section auto-updates the
% navigation dots shown in the Madrid header.
%
% For multi-paper talks, typical structure:
%   \section{Paper 1}
%   \section{Paper 2}
%   \section{Comparison}
%
% For figures, drop files into a `figures/` directory next to this .tex file
% and reference relatively: \includegraphics[width=...]{figures/myfig.pdf}
% ============================================================================

\documentclass[aspectratio=169]{beamer}

% --- Theme -------------------------------------------------------------------
\usetheme{Madrid}
\usecolortheme{default}

% CML navy blue — the shade used in the BigData_ML_Lecture reference
\definecolor{cmlblue}{RGB}{28, 53, 115}
% On Madrid, recoloring `structure` + the palette cascades cleanly to the
% title banner, the headline strip, the footer, and frame titles — without
% the title-text ending up the same color as its background. Don't override
% individual elements (title, frametitle) here unless you're willing to
% double-check that fg and bg haven't collided.
\setbeamercolor{structure}{fg=cmlblue}
\setbeamercolor{palette primary}{bg=cmlblue, fg=white}
\setbeamercolor{palette secondary}{bg=cmlblue!80, fg=white}
\setbeamercolor{palette tertiary}{bg=cmlblue!60, fg=white}

% --- Packages ---------------------------------------------------------------
\usepackage[T1]{fontenc}
\usepackage[utf8]{inputenc}
\usepackage{amsmath, amssymb, amsthm}
\usepackage{graphicx}
\usepackage{hyperref}
\hypersetup{colorlinks=true, linkcolor=cmlblue, urlcolor=cmlblue, citecolor=cmlblue!70}
\usepackage[round, authoryear]{natbib}
% Bibliography is loaded via \bibliography{} in the References frame near
% the end of the document. Natbib + bibtex works out of the box on any
% TeX Live install. Citation commands: \citet{key}, \citep{key}, \citeauthor{key}.

% --- Notes configuration ----------------------------------------------------
% Two compile modes:
%   default: talk-facing (notes hidden)
%   uncomment the show-only-notes line + recompile for a notes-only handout
\setbeameroption{hide notes}
% \setbeameroption{show only notes}

% --- Metadata ---------------------------------------------------------------
\title[How Memory Works: A Retrieved…]{How Memory Works: A Retrieved-Context Account}
\subtitle{Talk based on our 2024 manuscript}
\author[Alice Smith et al.]{Alice Smith, Bob Jones}
\institute[Penn CML]{Computational Memory Lab \\ University of Pennsylvania}
\date{\today}

% ============================================================================
\begin{document}

% ============================================================================
% Slides below generated by scripts/plan_deck.py
% ============================================================================

\begin{frame}[plain]
  \titlepage
\end{frame}

\begin{frame}{Outline}
  \tableofcontents
\end{frame}

\section{Motivation}

\begin{frame}{[TODO motivation slide 1]}
  \begin{itemize}
    \item {[TODO bullet 1]}
    \item {[TODO bullet 2]}
  \end{itemize}
  \note{[TODO speaker note]}
\end{frame}

\section{Background}

\begin{frame}{[TODO background 1]}
  \begin{itemize}
    \item {[TODO]}
  \end{itemize}
  \note{[TODO speaker note]}
\end{frame}

\section{Methods}

\begin{frame}{[TODO methods 1]}
  \begin{itemize}
    \item {[TODO design / participants / analyses]}
  \end{itemize}
  \note{[TODO]}
\end{frame}

\section{Results}

\begin{frame}{[TODO — figure 1]}
  \centering
  \includegraphics[width=0.85\textwidth,height=0.75\textheight,keepaspectratio]{figures/fig1.pdf}
  \\[0.5em]\footnotesize [TODO short caption]
  \note{[TODO walk the figure]}
\end{frame}

\begin{frame}{[TODO — figure 2]}
  \centering
  \includegraphics[width=0.85\textwidth,height=0.75\textheight,keepaspectratio]{figures/fig2.pdf}
  \\[0.5em]\footnotesize [TODO short caption]
  \note{[TODO walk the figure]}
\end{frame}

\begin{frame}{[TODO — figure 3]}
  \centering
  \includegraphics[width=0.85\textwidth,height=0.75\textheight,keepaspectratio]{figures/fig3.png}
  \\[0.5em]\footnotesize [TODO short caption]
  \note{[TODO walk the figure]}
\end{frame}

\section{Discussion}

\begin{frame}{[TODO implication 1]}
  \begin{itemize}
    \item {[TODO]}
  \end{itemize}
  \note{[TODO]}
\end{frame}

\begin{frame}{Acknowledgments}
  \begin{itemize}
    \item {[TODO collaborators + funding]}
  \end{itemize}
\end{frame}

\begin{frame}[allowframebreaks]{References}
  \bibliographystyle{apalike}
  \bibliography{references}
\end{frame}

\end{document}
